% !TeX root = ../../Book4.tex
% 纲要标题：### 13.3 乘积与括号
% 纲要标签：b4:sec:1203
% 纲要位置：计划纲要.md，第 3543 行。

\section{乘积与括号}
\label{b4:sec:1203}

余链值相乘与将余链代入另一个余链，是两种次数不同的运算。后者把结合律的高阶项写成括号，使延拓障碍成为具体的三次上同调类。

% 纲要内容（仅作写作计划，尚非教材正文）：
%
% * 杯积
% * Gerstenhaber 括号的构造思想
% * 形变障碍与 HH^3 的作用
%

\subsection{杯积}
\label{b4:subsec:1203:01}

当系数是 \(A\) 本身时，可以把两个上链的值相乘。另一种做法是把一个上链的值代入另一个上链。前者给出杯积，后者适合表达结合律及形变障碍；两种运算的次数不同。

\begin{definition}\label{b4:def:hochschild-cup}
设 \(k\) 为域，\(A\) 为结合含幺 \(k\) 代数，\(f\in C^p(A,A)\)、\(g\in C^q(A,A)\)，其中 \(p,q\ge0\)。规定
\[
(f\smile g)(a_1,\ldots,a_{p+q})
=f(a_1,\ldots,a_p)g(a_{p+1},\ldots,a_{p+q}).
\]
这个运算称为\term[beiji]{杯积}[cup product]。当某次为零时，相应上链就是一个 \(A\) 中元素，占据该因子的位置。
\end{definition}
\symidx{\(f\smile g\)：Hochschild 上链及上同调类的杯积}

\begin{proposition}\label{b4:prop:hochschild-cup-differential}
设 \(k\) 为域，\(A\) 为结合含幺 \(k\) 代数，\(p,q\ge0\)，\(f\in C^p(A,A)\)、\(g\in C^q(A,A)\)。\Cref{b4:def:hochschild-cup}定义的杯积结合，零次上链 \(1\) 是单位，且
\[
\delta(f\smile g)=(\delta f)\smile g+(-1)^p f\smile\delta g.
\]
因而它诱导 \(\operatorname{HH}^*(A)\) 上的带单位分次结合代数结构。
\end{proposition}
\begin{proof}
结合性来自 \(A\) 的乘法结合性。先看 \(p=q=1\) 的微分：
\[
\begin{aligned}
\delta(f\smile g)(a,b,c)
={}&a\cdot f(b)g(c)-f(ab)g(c)\\
&+f(a)g(bc)-f(a)g(b)\cdot c\\
={}&(\delta f)(a,b)g(c)-f(a)(\delta g)(b,c).
\end{aligned}
\]
右边展开时，\(f(a)b g(c)\) 的两次出现符号相反；这是一般次数中交界项相消的最小情形。在微分公式中，合并前 \(p+1\) 个位置产生 \((\delta f)\smile g\) 的各内部项，合并后 \(q+1\) 个位置产生 \((-1)^pf\smile\delta g\) 的各内部项。两部分交界处额外出现的项分别为
\((-1)^{p+1}f(a_1,\ldots,a_p)a_{p+1}\cdot g(\ldots)\) 与
\((-1)^pf(a_1,\ldots,a_p)a_{p+1}\cdot g(\ldots)\)，恰好抵消；最左、最右作用项保留下来，即为 \(\delta(f\smile g)\)。零次情形由同一左右乘法公式成立。

因此两个余圈的杯积仍是余圈。若 \(f=\delta u\)、\(\delta g=0\)，则 \(f\smile g=\delta(u\smile g)\)；另一因子为余边时同样成立。这证明乘法在上同调类上良定义。
\end{proof}

杯积也是 \(\operatorname{Ext}_{A^e}^*(A,A)\) 的 Yoneda 乘积。为直接比较两种定义，设 \(g\in C^q(A,A)\) 为余圈，在 bar 消解上定义
\[
\begin{aligned}
\widetilde g_i:B_{q+i}(A)&\longrightarrow B_i(A),\\
a_0\otimes\cdots\otimes a_{q+i+1}
&\longmapsto a_0\otimes\cdots\otimes a_i\otimes
g(a_{i+1},\ldots,a_{i+q})a_{q+i+1}\qquad(i\ge0).
\end{aligned}
\]
这是双模同态；\(\delta g=0\) 使合并尾部因子所得各项相消，因而 \(\mathrm{d}_i\widetilde g_i=\widetilde g_{i-1}\mathrm{d}_{q+i}\) 对 \(i\ge1\) 成立。增广与 \(\widetilde g_0\) 的复合就是代表 \(g\) 的 bar 余圈。若 \(f\in C^p(A,A)\) 也是余圈，把它与 \(\widetilde g_p\) 复合，在标准张量上便得到 \(f(a_1,\ldots,a_p)g(a_{p+1},\ldots,a_{p+q})\)。这正是\cref{b4:prop:yoneda-resolution-product}中用消解提升计算 Yoneda 乘积的公式。

\ProseParagraphBreak
例如中心中的元素通过零次杯积作用于所有上同调群。一般系数双模 \(M\) 没有自带的 \(M\otimes M\to M\)，所以不能不加数据就为 \(\operatorname{HH}^*(A,M)\) 写出同样的乘法。

\subsection{Gerstenhaber 括号的构造思想}
\label{b4:subsec:1203:02}

杯积把两组自变量前后排列，不能直接表达 \(\phi(\phi(a,b),c)\)。为比较这类嵌套乘法，需要记录插入的位置及其符号。例如两个二次上链的插入将是
\[
(f\circ g)(a,b,c)=f(g(a,b),c)-f(a,g(b,c)).
\]
一次插入把两个自变量合成一个值，所以结果有三个自变量。一般地，插入 \(q\) 次上链会给后续位置增加 \(q-1\) 个自变量；符号所使用的正是这个移后次数。

\begin{definition}\label{b4:def:gerstenhaber-bracket}
设 \(k\) 为域，\(A\) 为结合含幺 \(k\) 代数，\(p,q\ge0\)。对 \(f\in C^p(A,A)\)、\(g\in C^q(A,A)\)，在 \(p\ge1\) 时定义
\[
\begin{aligned}
f\circ g&=\sum_{i=0}^{p-1}(-1)^{i(q-1)}f\circ_i g,\\
(f\circ_i g)(a_1,\ldots,a_{p+q-1})
&=f(a_1,\ldots,a_i,g(a_{i+1},\ldots,a_{i+q}),
a_{i+q+1},\ldots,a_{p+q-1}).
\end{aligned}
\]
\(q=0\) 时将常数 \(g\in A\) 插入第 \(i+1\) 个位置；\(p=0\) 时规定 \(f\circ g=0\)。定义
\[
[f,g]=f\circ g-(-1)^{(p-1)(q-1)}g\circ f.
\]
这个次数为 \(-1\) 的运算称为 \term[gerstenhaber-kuohao]{Gerstenhaber 括号}[Gerstenhaber bracket]。\footnote{格斯滕哈伯（Murray Gerstenhaber，1927-2024），美国数学家，研究代数形变理论，并引入以他命名的括号与代数结构。}两个零次上链的括号取零。
\end{definition}
\symidx{\([f,g]\)：Hochschild 上链上的 Gerstenhaber 括号}

零次上链也参与插入。例如 \(D\in C^1(A,A)\)、\(a\in C^0(A,A)=A\) 满足 \(D\circ a=D(a)\)、\(a\circ D=0\)，故 \([D,a]=D(a)\)。它是一个零次上链，次数仍为 \(1+0-1\)。

\begin{proposition}\label{b4:prop:gerstenhaber-identities}
设 \(k\) 为域，\(A\) 为结合含幺 \(k\) 代数，\(p,q\ge0\)，以 \(|f|=p-1\) 表示 \(f\in C^p(A,A)\) 的移后次数。\Cref{b4:def:gerstenhaber-bracket}定义的括号满足分次反对称性及分次 Jacobi 恒等式，且诱导
\[
\operatorname{HH}^p(A)\times\operatorname{HH}^q(A)
\longrightarrow\operatorname{HH}^{p+q-1}(A).
\]
这里约定负次数的上同调为零。若 \(m(a,b)=ab\) 为原乘法，则
\[
[m,f]=(-1)^{p-1}\delta f.
\]
此外，杯积在上同调上满足分次交换律
\[
\alpha\smile\beta=(-1)^{pq}\beta\smile\alpha,
\qquad \alpha\in\operatorname{HH}^p(A),\quad
\beta\in\operatorname{HH}^q(A).
\]
\end{proposition}
\begin{proof}
记 \(a=|f|\)、\(b=|g|\)、\(c=|h|\)，并令
\[
P(f,g,h)=(f\circ g)\circ h-f\circ(g\circ h).
\]
先令 \(f,g,h\) 都为二次上链，观察两种实际复合。左边的树把 \(h\) 插入 \(g\) 内部，再把 \(g\) 插入 \(f\)；右边把 \(g,h\) 插入 \(f\) 的两个不同位置。箭头表示把下方运算的输出代入上方运算。
\[
\begin{tikzpicture}
\node (nf) at (-3,0) {\(f\)};
\node (ng) at (-3.7,-0.9) {\(g\)};
\node (na4) at (-2.2,-0.9) {\(a_4\)};
\node (nh) at (-4.2,-1.8) {\(h\)};
\node (na3) at (-3.1,-1.8) {\(a_3\)};
\node (na1) at (-4.7,-2.7) {\(a_1\)};
\node (na2) at (-3.7,-2.7) {\(a_2\)};
\draw[->] (ng)--(nf); \draw[->] (na4)--(nf);
\draw[->] (nh)--(ng); \draw[->] (na3)--(ng);
\draw[->] (na1)--(nh); \draw[->] (na2)--(nh);
\node (df) at (2.3,0) {\(f\)};
\node (dg) at (1.4,-0.9) {\(g\)};
\node (dh) at (3.2,-0.9) {\(h\)};
\node (da1) at (0.9,-1.8) {\(a_1\)};
\node (da2) at (1.9,-1.8) {\(a_2\)};
\node (da3) at (2.7,-1.8) {\(a_3\)};
\node (da4) at (3.7,-1.8) {\(a_4\)};
\draw[->] (dg)--(df); \draw[->] (dh)--(df);
\draw[->] (da1)--(dg); \draw[->] (da2)--(dg);
\draw[->] (da3)--(dh); \draw[->] (da4)--(dh);
\end{tikzpicture}
\]
左边的 \(f(g(h(a_1,a_2),a_3),a_4)\) 同时出现在 \((f\circ g)\circ h\) 与 \(f\circ(g\circ h)\) 中，符号均为正，所以在 \(P\) 中抵消；把 \(h\) 插入 \(g\) 的另一个位置时，两处符号均为负，也抵消。右边的 \(f(g(a_1,a_2),h(a_3,a_4))\) 则不属于 \(f\circ(g\circ h)\)。在 \((f\circ g)\circ h\) 中，先占 \(f\) 的位置 \(0\) 后，\(h\) 应插入新位置 \(2\)，符号为 \((-1)^0(-1)^2=1\)；若先在 \(f\) 的位置 \(1\) 插入 \(h\)，再在位置 \(0\) 插入 \(g\)，符号为 \((-1)^1(-1)^0=-1\)。因此这一次展开后只剩
\[
\begin{aligned}
P(f,g,h)(a_1,a_2,a_3,a_4)
={}&f(g(a_1,a_2),h(a_3,a_4))\\
&-f(h(a_1,a_2),g(a_3,a_4)).
\end{aligned}
\]
交换 \(g,h\) 恰使这个表达式变号。一般次数的计算整理的是同样的嵌套抵消与异位配对。插入后的自变量数为 \(p+q-1\)，减去一后为 \((p-1)+(q-1)\)，所以移后次数恰使插入成为次数相加的运算。

以下 \(i,j\) 表示从零开始计数的插入位置，而 \(a,b,c\) 表示移后次数；它们与图中的自变量 \(a_1,\ldots,a_4\) 无关。在 \(f\) 的位置 \(i\) 插入 \(g\)，再在 \(g\) 的位置 \(j\) 插入 \(h\)，两种嵌套次序的符号分别为
\[
(-1)^{ib+(i+j)c}
=(-1)^{i(b+c)+jc}.
\]
因此这类内部插入项在 \(P(f,g,h)\) 中抵消，只剩把 \(g,h\) 放入 \(f\) 的不同位置的项。

若 \(g\) 占原位置 \(i\)，\(h\) 占原位置 \(j\)，且 \(i<j\)，先插 \(g\) 后插 \(h\) 时，\(h\) 的位置变为 \(j+b\)，符号为
\[
(-1)^{ib+(j+b)c};
\]
先插 \(h\) 后插 \(g\) 的符号为 \((-1)^{jc+ib}\)，两者之比是 \((-1)^{bc}\)。当 \(j<i\) 时，两种符号分别为
\((-1)^{ib+jc}\) 与 \((-1)^{jc+(i+c)b}\)，仍相差 \((-1)^{bc}\)。于是
\[
P(f,g,h)=(-1)^{bc}P(f,h,g).
\]
这就是带符号的右 pre-Lie 恒等式。若 \(g\) 或 \(h\) 为零次上链，内部位置为空，插入常数则使后续位置前移一格，即相应移后次数为 \(-1\)；上面的异位配对仍成立。若 \(f\) 为零次上链，两边均为零。所有落入负上链次数的表达式也取零。

括号的分次反对称性直接来自定义。为核验 Jacobi 恒等式，把其导子形式的差展开为
\[
\begin{aligned}
&[f,[g,h]]-[[f,g],h]-(-1)^{ab}[g,[f,h]]\\
={}&-P(f,g,h)+(-1)^{bc}P(f,h,g)\\
&+(-1)^{ab}\bigl(P(g,f,h)-(-1)^{ac}P(g,h,f)\bigr)\\
&-(-1)^{c(a+b)}\bigl(P(h,f,g)-(-1)^{ab}P(h,g,f)\bigr).
\end{aligned}
\]
第一对由 \(P(f,g,h)=(-1)^{bc}P(f,h,g)\) 抵消；后两对分别交换 \(f,h\) 和 \(f,g\)，由同一恒等式抵消，故差为零。再用分次反对称性改写，得到通常的循环形式
\[
(-1)^{|f|\cdot|h|}[f,[g,h]]
+(-1)^{|g|\cdot|f|}[g,[h,f]]
+(-1)^{|h|\cdot|g|}[h,[f,g]]=0.
\]
将二次上链 \(m\) 代入括号定义，最外两项是左右乘法，内部各项是合并相邻自变量，给出 \([m,f]=(-1)^{p-1}\delta f\)。例如对二次上链 \(\phi\)，
\[
\begin{aligned}
[m,\phi](a,b,c)
={}&\phi(a,b)\cdot c-a\cdot\phi(b,c)\\
&+\phi(ab,c)-\phi(a,bc)
=-(\delta\phi)(a,b,c).
\end{aligned}
\]
原乘法的结合律就是 \(m\circ m=0\)，而一阶结合律就是 \([m,\phi]=0\)。特别地，令 \(D f=[m,f]\)，Jacobi 恒等式\footnote{雅可比（Carl Gustav Jacob Jacobi，1804-1851），德国数学家，在椭圆函数和动力学的微分方程方面作出重要贡献。}给出
\[
D[f,g]=[Df,g]+(-1)^{|f|}[f,Dg].
\]
\(D\) 与 \(\delta\) 在每次只差一个符号，因此有相同的余圈、余边。最后一个等式说明余圈的括号仍为余圈，且一个因子为余边、另一个为余圈时结果为余边。这就使括号下降到上同调。

最后把原乘法放在异位插入公式的首位。由于 \(m\) 只有两个位置，其余两项给出
\[
P(m,f,g)=(-1)^{p(q-1)}f\smile g+(-1)^{p-1}g\smile f.
\]
将 \([m,f\circ g]\) 展开，并用
\(P(f,m,g)=(-1)^{q-1}P(f,g,m)\)，得到
\[
\begin{aligned}
\delta(f\circ g)
={}&(-1)^{q-1}(\delta f)\circ g+f\circ\delta g\\
&+(-1)^{q-1}\bigl((-1)^{pq}f\smile g-g\smile f\bigr).
\end{aligned}
\]
若 \(f,g\) 都是余圈，右边最后一行就是余边，故上同调上的杯积满足分次交换律。例如对导子 \(D,E\)，这个等式成为
\[
\delta(D\circ E)=-(D\smile E+E\smile D).
\]
两个杯积在上链层面未必互为负号，其和却是明确的余边，所以上同调类互为负号。零次情形同样包含在一般公式中；当两者都为零次余圈时，它们是中心元素，杯积直接交换。
\end{proof}

杯积与括号还满足分次 Leibniz 律：对
\(\alpha\in\operatorname{HH}^p(A)\)、\(\beta\in\operatorname{HH}^q(A)\)、\(\gamma\in\operatorname{HH}^r(A)\)，有
\[
[\alpha\smile\beta,\gamma]
=[\alpha,\gamma]\smile\beta
+(-1)^{p(r-1)}\alpha\smile[\beta,\gamma].
\]
按本节的括号符号约定，固定右变量 \(\gamma\) 得到杯积的次数为 \(r-1\) 的分次导子。这使 \(\operatorname{HH}^*(A)\) 同时具有相容的分次交换乘法与次数为 \(-1\) 的 Lie 括号，合称 Gerstenhaber 代数结构。\cite[Section 8, Theorem 5, Corollaries 1--2, pp. 285--287]{Gerstenhaber1963CohomologyStructure}
\begin{proofsketch}
取代表 \(\alpha,\beta,\gamma\) 的余圈 \(f,g,h\)。将 \(h\) 插入杯积的前半或后半，直接得到
\[
(f\smile g)\circ h
=(f\circ h)\smile g+(-1)^{p(r-1)}f\smile(g\circ h).
\]
另一方向需比较 \(h\circ(f\smile g)\) 与
\((-1)^{q(r-1)}(h\circ f)\smile g+f\smile(h\circ g)\)。所用的双重插入上链可以写成
\[
H=\sum_{i=0}^{r-2}\ \sum_{j=i+p}^{p+r-2}
(-1)^{(p-1)i+(q-1)j}(h\circ_i f)\circ_j g.
\]
先把 \(f\) 插入 \(h\) 的位置 \(i\)，再把 \(g\) 插入其右方的另一个原位置；前次插入增加了 \(p-1\) 个位置，故第二个指标从 \(i+p\) 开始。\(H\) 的次数为 \(p+q+r-2\)，在 \(r<2\) 时取空和。展开 \(\delta H\) 时，内部合并项由 \(\delta f=\delta g=\delta h=0\) 抵消，异位项按位置配对，交界项给出上述比较之差的一个整体符号倍数。因此该差为余边。结合杯积插入等式与括号定义，便得到所需 Leibniz 律。这里省略 \(\delta H\) 的全部位置符号展开；完整计算见上述文献。
\end{proofsketch}

对一次上链，即线性自映射，括号为普通交换子 \([D,E]=D\circ E-E\circ D\)。因此导子的交换子仍是导子，并给出 \(\operatorname{HH}^1(A)\) 上的 Lie 括号\footnote{李（Marius Sophus Lie，1842-1899），挪威数学家，创立连续变换群理论，奠定 Lie 群与 Lie 代数的基础。}。对二次上链 \(\phi\)，则
\[
(\phi\circ\phi)(a,b,c)
=\phi(\phi(a,b),c)-\phi(a,\phi(b,c)).
\]
当 \(\operatorname{char}k\ne2\) 时，它等于 \(\dfrac12[\phi,\phi]\)。在特征二中不能用除以二的写法；实际障碍始终是左边这个插入表达式。

\subsection{\texorpdfstring{形变障碍与 \(\operatorname{HH}^3\)}{形变障碍与 HH3}}
\label{b4:subsec:1203:03}

一次结合律只给线性方程，第二次就出现刚才的二次表达式。所有阶的形变需要先确定形式级数的含义。设 \(k\) 为域，\(A\) 为结合含幺 \(k\) 代数，记
\[
A[\![t]\!]
=\varprojlim_{N\ge0}\bigl(A\otimes_k k[t]/(t^{N+1})\bigr)
\cong\prod_{n\ge0}At^n.
\]
其元素为 \(\sum_{n\ge0}a_nt^n\)，\(t\) 进拓扑以 \(t^{N+1}A[\![t]\!]\) 为零点邻域基。若 \(A\) 无限维，这一完备化一般不同于普通张量积 \(A\otimes_k k[\![t]\!]\)。

\ProseParagraphBreak

例如取 \(A=k[x]\)，则 \(\sum_{n\ge0}x^nt^n\) 是 \(A[\![t]\!]\) 中的元素。普通张量积中的元素却只能是有限和 \(\sum_{i=1}^r a_i\otimes f_i(t)\)。把它展开为 \(t\) 的级数后，每个系数都落在同一个有限维空间 \(\operatorname{span}_k\{a_1,\ldots,a_r\}\) 中；而 \(1,x,x^2,\ldots\) 线性无关，不可能全部落在这种空间里。因此该级数不来自普通张量积，完备化允许的是这样的无限系数族。

\ProseParagraphBreak
在这一完备模上，保持单位的形式形变是连续的 \(k[\![t]\!]\) 双线性结合乘法，模 \(t\) 后等于 \(A\) 的原乘法，单位仍为 \(1\)。它在常数元素上的值唯一写成
\[
a*b=ab+\sum_{i\ge1}t^i\mu_i(a,b),
\]
其中 \(\mu_i\) 为归一化 \(k\) 双线性映射。两个这样的乘法若由保持单位、模 \(t\) 为恒等的连续 \(k[\![t]\!]\) 线性同构 \(T\) 相联系，即
\(T(a*b)=T(a)*'T(b)\)，就称为等价。这样的坐标变换写成
\(T=\operatorname{id}+\sum_{i\ge1}t^ih_i\)，其中 \(h_i:A\to A\) 为 \(k\) 线性映射，且 \(h_i(1)=0\)；其逆可在各有限截断中逐次求出。

\ProseParagraphBreak
令 \(\mu_0=m\) 为原乘法。双线性与连续性使一般级数的乘积必须为
\[
\left(\sum_{r\ge0}a_rt^r\right)*
\left(\sum_{s\ge0}b_st^s\right)
=\sum_{n\ge0}t^n
\sum_{\substack{i+r+s=n\\i,r,s\ge0}}\mu_i(a_r,b_s).
\]
固定 \(n\) 时，内层只取满足 \(i+r+s=n\) 的有限多个三元组。因此每个系数都已确定，模 \(t^{N+1}\) 的乘积只依赖两个输入的相应截断；这也说明该公式给出连续的 \(k[\![t]\!]\) 双线性乘法。

\ProseParagraphBreak
若已知 \(\phi\) 是归一化 \(2\)-余圈，尝试令
\[
a*b=ab+t\cdot\phi(a,b)+t^2\cdot\psi(a,b)\pmod {t^3}.
\]
比较结合律的二次系数，得到
\[
\begin{aligned}
\psi(ab,c)+\psi(a,b)\cdot c+\phi(\phi(a,b),c)
={}&\psi(a,bc)+a\cdot\psi(b,c)\\
&+\phi(a,\phi(b,c)).
\end{aligned}
\]
前四个 \(\psi\) 项来自二阶修正与原乘法的结合，两个嵌套 \(\phi\) 项则来自两次一阶修正。将 \(\psi\) 项移到一侧，所需且仅需的条件是
\begin{equation}\label{b4:eq:second-deformation-obstruction}
\delta\psi=\phi\circ\phi.
\end{equation}
因此出现一个三次上同调问题：右边不仅要是余圈，还必须是余边。

\begin{theorem}\label{b4:thm:hochschild-obstruction}
设 \(k\) 为域，\(A\) 为非零结合含幺 \(k\) 代数，\(m(a,b)=ab\) 为其乘法，\(N\ge1\)，\(D_N=k[t]/(t^{N+1})\)。在 \(A\otimes_kD_N\) 上给定保持单位的结合 \(D_N\) 双线性乘法
\[
m_t=m+t\mu_1+\cdots+t^N\mu_N,
\]
其中各 \(\mu_i:A\otimes_kA\to A\) 为 \(k\) 线性映射，满足 \(\mu_i(1,a)=\mu_i(a,1)=0\) 对所有 \(a\in A\) 成立。令
\[
O_{N+1}(a,b,c)=
\sum_{\substack{i+j=N+1\\i,j\ge1}}
\bigl(\mu_i(\mu_j(a,b),c)-\mu_i(a,\mu_j(b,c))\bigr).
\]
则 \(O_{N+1}\) 是归一化三余圈。该给定形变能延拓到模 \(t^{N+2}\) 的保持单位形变，当且仅当
\([O_{N+1}]=0\) 于 \(\operatorname{HH}^3(A)\)。
\end{theorem}
\begin{proof}
在现有截断多项式乘法后加上 \(t^{N+1}\nu\)，结合子的 \(t^{N+1}\) 系数为 \(O_{N+1}-\delta\nu\)，因为涉及原乘法 \(m\) 与 \(\nu\) 的四项恰好是 \(-\delta\nu\)。故延拓要求 \(\delta\nu=O_{N+1}\)。

为证明右边是余圈，不需要除以任何整数。对任意双线性乘法 \(*\)，记其结合子为 \(F(a,b,c)=(a*b)*c-a*(b*c)\)。直接展开可得
\[
\begin{aligned}
0={}&a*F(b,c,d)-F(a*b,c,d)+F(a,b*c,d)\\
&-F(a,b,c*d)+F(a,b,c)*d.
\end{aligned}
\]
展开后出现的五种四重乘积是
\[
\begin{gathered}
((a*b)*c)*d,\qquad (a*(b*c))*d,\qquad (a*b)*(c*d),\\
a*((b*c)*d),\qquad a*(b*(c*d)).
\end{gathered}
\]
每一种恰出现两次，符号相反；这一步未使用 \(*\) 的结合性。对尚未加入 \(\nu\) 的 \(m_t\) 使用此式。因结合子的所有低于 \(N+1\) 次系数为零，比较 \(t^{N+1}\) 系数时，外部乘法只取原乘法，所得正是 \(\delta O_{N+1}=0\)。各 \(\mu_i\) 归一化也使 \(O_{N+1}\) 任一自变量为 \(1\) 时为零。

若延拓存在，\(O_{N+1}=\delta\nu\) 已说明其类为零。反过来，若类为零，先取任意 \(\nu\) 使 \(\delta\nu=O_{N+1}\)。右边在带单位的三元组上为零，所以与\cref{b4:thm:hh2-deformation-classification}中相同的两次代入给出
\(\nu(1,a)=ca\)、\(\nu(a,1)=ac\)，其中 \(c=\nu(1,1)\)。选择 \(h(1)=c\)，以 \(\nu-\delta h\) 替代 \(\nu\)，即可既保持 \(\delta\nu=O_{N+1}\)，又使 \(\nu\) 归一化。所得延拓保持单位。
\end{proof}

\begin{corollary}\label{b4:cor:hh3-unobstructed}
设 \(k\) 为域，\(A\) 为非零结合含幺 \(k\) 代数，\(m(a,b)=ab\) 为其乘法。若 \(\operatorname{HH}^3(A)=0\)，则每个保持单位的一阶形变都能逐次延拓成 \(A[\![t]\!]\) 上的结合形式乘法
\(m_t=m+\sum_{i\ge1}t^i\mu_i\)，其中各 \(\mu_i:A\otimes_kA\to A\) 为归一化 \(k\) 线性映射。
\end{corollary}
\begin{proof}
从给定的一阶形变 \(m^{(1)}\) 开始，用\cref{b4:thm:hochschild-obstruction}归纳选择 \(\mu_{N+1}\)，得到保持单位的结合截断乘法 \(m^{(N+1)}\)，并使其模 \(t^{N+1}\) 恰为已选的 \(m^{(N)}\)。因此这些乘法构成兼容系统，可取逆极限。所得级数乘法就是前面的一般系数公式；它连续且 \(k[\![t]\!]\) 双线性。每个次数的结合律已在相应截断中成立，单位由所有 \(\mu_i\) 的归一化保证。
\end{proof}

这是形式延拓，不涉及实数或复数分析中的收敛。若 \(\operatorname{HH}^3(A)\ne0\)，也不能反向断言所有形变都会受阻；要计算的是当前已选系数产生的特定类 \([O_{N+1}]\)。后续系数的不同选择还可能影响更高阶障碍。下面的代数说明，三次上同调非零也可能与所有一阶形变都可延拓同时发生。

\begin{proposition}\label{b4:prop:dual-numbers-unobstructed}
设 \(k\) 为任意域，\(A=k[x]/(x^2)\)。则 \(\operatorname{HH}^3(A)\ne0\)，但 \(A\) 的每个保持单位的一阶形变都能延拓成 \(A[\![t]\!]\) 上的保持单位形式形变。
\end{proposition}
\begin{proof}
定义三线性上链 \(\eta\)，使 \(\eta(x,x,x)=x\)，且在基 \(1,x\) 组成的三元组中，只要有一个自变量为 \(1\)，其值就为零。归一化使 \(\delta\eta\) 在带单位的四元组上为零；在唯一剩下的基四元组上，
\[
(\delta\eta)(x,x,x,x)
=x\eta(x,x,x)+\eta(x,x,x)x=2x^2=0.
\]
因此 \(\eta\) 是余圈。另一方面，对任意二次上链 \(\psi\)，由 \(x^2=0\) 和交换性得
\[
(\delta\psi)(x,x,x)=x\psi(x,x)-\psi(x,x)x=0.
\]
它不可能等于取值为 \(x\ne0\) 的 \(\eta\)，故 \([\eta]\ne0\)。

任一保持单位的一阶形变由归一化余圈 \(\phi(x,x)=u+vx\)、\(u,v\in k\) 确定。取结合代数
\[
k[\![t]\!][X]/\bigl(X^2-t(u+vX)\bigr).
\]
首一多项式除法使它在 \(k[\![t]\!]\) 上以 \(1,X\) 为自由基，故可与 \(A[\![t]\!]\) 识别；其乘法连续，单位保持不变，模 \(t^2\) 后恰为给定的一阶形变。这在任何特征下都成立。
\end{proof}

\ProseParagraphBreak
低阶群的含义可以合在一起比较：零次记录中心化与交换子商，一次记录导子与微分，二次分类一阶乘法形变，三次容纳逐次延拓的障碍类。具体对象、系数和条件见\cref{b4:tab:hochschild-low-degrees}。
\begin{table}[htbp]
\centering
\caption[Hochschild 理论的低次数解释]{Hochschild 理论的低次数解释。这里 \(k\) 为域，\(A\) 为非零结合含幺 \(k\) 代数，\(M\) 为 \(A\) 双模。}
\label{b4:tab:hochschild-low-degrees}
\begin{tabularx}{\textwidth}{@{}p{0.22\textwidth}>{\raggedright\arraybackslash}X@{}}
\toprule
对象 & 代数意义 \\
\midrule
\(\operatorname{HH}_0(A,M)\)
& \(M/[A,M]\)，消去左右作用之差 \\
\(\operatorname{HH}^0(A,M)\)
& \(\{\,m\in M\mid am=ma\;\text{对所有}\;a\in A\,\}\) \\
\(\operatorname{HH}_1(A)\)
& 当 \(A\) 交换时，等于 Kähler 微分模 \(\Omega_{A/k}\) \\
\(\operatorname{HH}^1(A,M)\)
& \(\operatorname{Der}_k(A,M)/\operatorname{Inn}_k(A,M)\) \\
\(\operatorname{HH}^2(A)\)
& 保持单位及固定还原识别的一阶乘法形变等价类 \\
\(\operatorname{HH}^2(A,M)\)
& 固定核 \(M\)、商 \(A\) 及双模作用的平方零代数扩张等价类 \\
\(\operatorname{HH}^3(A)\)
& 给定有限阶形变的延拓障碍类所在的空间；并非每个三次类都必须出现为障碍 \\
\bottomrule
\end{tabularx}
\end{table}

